\chapter{Conclusiones}
\label{cap:conclusiones}

\begin{pendiente}[Por completar / personalizar]
Las conclusiones que siguen son una propuesta redactada a partir del trabajo efectivamente
realizado. Conviene ajustarlas a la experiencia concreta del equipo y agregar el aprendizaje
personal de cada integrante, que es lo que el profesor suele valorar en esta sección.
\end{pendiente}

\textbf{Sobre el modelado del sistema.} El ejercicio confirma que es posible reconstruir la
arquitectura funcional y estructural de un sistema en producción \textbf{sin acceso a su código
fuente}, partiendo de la descripción del proceso de negocio y de la observación de su
comportamiento externo. El modelo UML resultante no pretende ser idéntico a la implementación real
de BeeTracer, pero sí constituye un \textbf{plano coherente y completo} que explica todo el
comportamiento observado y que serviría de base para construir o extender un sistema equivalente.

\textbf{Sobre el traslado del modelo estructurado al orientado a objetos.} El paso del Informe 1
---con sus diagramas de flujo de datos y su modelo entidad-relación--- a este informe evidenció una
diferencia importante de enfoque. Mientras el modelo estructurado describe \textbf{qué
transformaciones sufre la información}, el modelo orientado a objetos obliga a decidir
\textbf{quién es responsable de cada comportamiento}. Esa decisión no es trivial y en varios puntos
del diseño fue necesario resolverla explícitamente, como al separar \clase{SincronizadorOffline}
del \clase{GestorFaena} o al escoger entre asociar las fotografías a la etapa o a la incidencia.

\textbf{Sobre el dominio.} El estudio del caso reveló que el valor central de BeeTracer \textbf{no
es tecnológico sino probatorio}. El sistema no compite por procesar más rápido ni por almacenar más
datos: compite por producir, en el momento exacto en que la carga cambia de custodia, una evidencia
fechada, geolocalizada y difícil de refutar. Toda su arquitectura se subordina a ese objetivo: por
eso la aplicación \textbf{bloquea el avance} si falta una fotografía, por eso sube las imágenes una
por una en lugar de esperar al final, y por eso existe una etapa de revisión humana antes de que el
reporte llegue al cliente.

\textbf{Sobre las restricciones del entorno como motor del diseño.} Uno de los aprendizajes más
claros del caso es que \textbf{las decisiones de arquitectura más determinantes no nacieron de
requisitos funcionales, sino de restricciones físicas del entorno}. El bloqueo de la señal Wi-Fi
por los contenedores metálicos explica el envío incremental, la compresión en dispositivo y el modo
offline completo. El hecho de que el operario trabaje con guantes y bajo presión de tiempo explica
el flujo guiado paso a paso. La sensibilidad comercial de la evidencia explica la eliminación de la
visualización en vivo por parte del cliente. Un modelo construido solo desde los requisitos
declarados habría omitido lo más característico del sistema.

\textbf{Sobre las limitaciones del trabajo.} El modelo presentado tiene un grado de inferencia
significativo. Las clases, atributos y multiplicidades son una propuesta coherente, pero no
verificada contra la implementación real. Los módulos complementarios ---visor 3D y aplicación de
avisos de transporte--- quedaron fuera del alcance por falta de información. Asimismo, no fue
posible validar el modelo con el equipo desarrollador, lo que habría permitido contrastar las
hipótesis de diseño.

\textbf{Proyección.} El modelo obtenido habilita al menos tres líneas de trabajo futuro:
\textbf{(i)} la incorporación de grabación de video del desconsolidado, funcionalidad que el propio
equipo de BeeTracer evalúa y que el modelo admitiría generalizando \clase{Fotografia} hacia una
clase \clase{Evidencia}; \textbf{(ii)} la integración del visor 3D de patios, añadiendo coordenadas
de ubicación a \clase{Contenedor}; y \textbf{(iii)} la explotación analítica del historial de
incidencias por cliente y por naviera, hoy registrado pero no aprovechado como indicador de
gestión.

\begin{pendiente}
Agregar aquí las conclusiones personales del equipo: dificultades encontradas durante el modelado,
decisiones que se discutieron y cómo se resolvieron, y aprendizaje obtenido de la asignatura.
\end{pendiente}
